% /root/monodromy, 2026-11-19. Proved; independent root, reduction_descent and obstructions_crystals proof audits passed.
\subsection{All invariant covectors on the supersingular product}
\label{subsec:ss-exact-covectors}
\status{Proved and checked. The parent and two other agents independently
audited the complete proof; the local Cartier certificate was independently
rerun by the parent. The persistent coherent nilpotent part is handled by a
bounded Frobenius-root extension and a homogeneous-coefficient argument.}

\begin{theorem}\label{thm:ss-product-exact-covectors}
Let $E_0:y^2+y=x^3$ over a perfect field of characteristic two, let
$X=E_0\times E_0$ and
$L=\mathcal O_{E_0}(3O)\boxtimes\mathcal O_{E_0}(3O)$.
For every $d\ge1$, the entire geometric and arithmetic invariant
covector spaces of the relative Hasse--Witt local system have dimension
\[
 \dim(V_d^\vee)^{G_{\rm geom}}
 =\dim(V_d^\vee)^{G_{\rm arith}}
 =\begin{cases}0,&d\text{ odd},\\1,&d\text{ even}.
 \end{cases}
\]
In even degree the unique nonzero covector is represented under
adjunction by $D(s)\omega$, where
$\omega=dx_1\wedge dx_2$, $D=\delta_1\delta_2$, and $\delta_i$
is the invariant derivation dual to $dx_i$, lifted to $L^d$ by its
canonical Frobenius-pullback connection. Thus every even-degree image
is contained in the pointwise stabilizer of one nonzero covector.
This theorem determines the entire invariant quotient, not the full group.
\end{theorem}

\begin{proof}
First work over an algebraic closure. Put $m=3d$,
$N=\mathcal O(mO)\boxtimes\mathcal O(mO)$,
$W=H^0(X,N)$ and $b=m^2$. On $P=\mathbf P(W)$ the full coherent
bundle is $\mathcal B=W^*\otimes\mathcal O_P(-1)$, with semilinear
Frobenius matrix $M(s)^t$. Here $M(s)$ is linear in the coefficient
vector $s$ and is the root-removed Cartier multiplication matrix.
Its relative subbundle is $\mathcal E=\Omega_P^1$; the quotient
$\mathcal B/\mathcal E=H^2(X,\mathcal O_X)\otimes\mathcal O_P$
has zero Frobenius. Thus the unit Lang local systems of the full
and relative bundles coincide.

\paragraph{The generic linearized image is already the stable image.}
The identity $C(\omega)=0$ gives $M(s)s=0$ in every degree.
For odd $m$, Theorem~\ref{thm:DE-A3-ss-product-rank} gives stable
rank $b-1$, so the generic linearized rank is also $b-1$.
For even $m$, $N$ is the relative Frobenius pullback of $N_{m/2}$.
The invariant derivations act on its sections, since they kill the
square transition functions of pullback frames. They commute and
satisfy $\delta_1^2=\delta_2^2=0$. A local $2$-basis calculation gives
\[
 C(s\delta_1(s)\omega)=C(s\delta_2(s)\omega)=0.
\]
Thus $s,\delta_1(s),\delta_2(s)$ lie in $\ker M(s)$.
They are generically independent: for the section $s=x_1x_2$
they are $x_1x_2,x_2,x_1$. Hence the generic linearized rank is
at most $b-3$. The exact stable-rank theorem gives $b-3$, so
equality holds. In both parities the image of the first linearized
Frobenius at the generic point is therefore its unit subspace.

\paragraph{A purely inseparable lift of an arbitrary invariant covector.}
Let $K=k(P)$ and let $I\subset\mathcal B_K$ be that unit image.
The restriction $F^*I\to I$ is invertible. The finite Lang-frame
equivalence identifies a geometric invariant covector with a
$K$-linear map $f:I\to K$ satisfying
$f(Fx)=f(x)^2$ on $I$. Define a functional on the full generic
bundle, with coefficients in $K^{1/2}$, by
\begin{equation}\label{eq:ss-covector-root-extension}
                  \lambda(x)^2=f(Fx).
\end{equation}
The right side is linear in the squares of the coordinates of $x$,
so this defines a linear functional over $K^{1/2}$. On $I$ it
equals $f$ by uniqueness of square roots. Since $Fx$ lies in $I$,
it satisfies $\lambda(Fx)=\lambda(x)^2$ on the whole bundle.
This extension is unique: any equivariant functional zero on $I$
has square zero by \eqref{eq:ss-covector-root-extension}.

In coordinates, $M(s)\lambda=\lambda^{[2]}$. The squared coefficient
vector $w=\lambda^{[2]}$ is $K$-rational and is a rational section of
$F_P^*\mathcal B^*=W\otimes\mathcal O_P(2)$. It satisfies
\begin{equation}\label{eq:ss-square-covector-equation}
                       M(s)^{[2]}w=w^{[2]}.
\end{equation}
It has no divisorial poles. In a regular local frame, if the largest
pole order among its coordinates were $a>0$, the right side would
have a coordinate of pole order $2a$, whereas the left side has
pole order at most $a$. Normality and local freeness extend the
section also across codimension two. Thus $w$ is a vector of
homogeneous quadratic polynomials in the coefficients $s_i$.

\paragraph{Cross terms cannot occur.}
Write $M(s)=\sum s_i M_i$. The common kernel of the $M_i$ is zero.
Indeed, relative finite-Frobenius duality, with the invariant top
form trivializing the canonical bundle, gives a perfect symmetric
Cartier pairing
\begin{equation}\label{eq:ss-perfect-Cartier-pairing}
 F_{X/k*}N\otimes F_{X/k*}N\longrightarrow N^{(2)}.
\end{equation}
The finite-duality and coefficient-trace formulas used here are
\cite[Lemma~1.3.6, Proposition~1.3.7, Theorem~1.3.8, pp.~26--28]
{BrionKumar2005}, already used in the general adjunction proof.
Its pairing on global sections is the root-removed bilinear map
$\beta(s,t)=M(s)t$. The bundle $F_{X/k*}N$ is globally generated.
On either elliptic factor, evaluation on a Frobenius fibre is
evaluation on a length-two subscheme, and the obstruction group is
$H^1(\mathcal O(mO)(-Z))=0$ because $m-2>0$.
Tensoring the two evaluation surjections proves the assertion on
the product. A global section pairing to zero with every global
section therefore vanishes in every fibre, by perfectness, and
is zero. This proves the common-kernel assertion, also after
squaring the matrices.

For $j<k$, let $c_{jk}$ be the coefficient of $s_js_k$ in $w$.
The right side of \eqref{eq:ss-square-covector-equation} has only
even exponents. The monomials with precisely $s_j,s_k$ occurring
to odd powers on its left side give
\[
              \sum_i s_i^2M_i^{[2]}c_{jk}=0.
\]
Thus every $M_i^{[2]}c_{jk}=0$, and $c_{jk}=0$.
All cross terms vanish. Since $k$ is perfect,
\[
                         w=D(s)^{[2]}
\]
for a $k$-linear endomorphism $D$ of $W$. Taking square roots of
\eqref{eq:ss-square-covector-equation} shows
\begin{equation}\label{eq:ss-exact-linear-operator}
                 \beta(s,D(s))=D(s)^{[2]}.
\end{equation}
The original functional is exactly $D(s)$, not merely its
purely inseparable extension. Conversely every $D$ satisfying
\eqref{eq:ss-exact-linear-operator} is a regular equivariant
covector. No quotient by scalar identities is taken here:
the full-to-relative quotient has Frobenius zero, so it has no
nonzero equivariant covectors.

\paragraph{Descent to a small endomorphism algebra.}
Polarizing \eqref{eq:ss-exact-linear-operator} gives
\[
                   \beta(Ds,t)=\beta(s,Dt).
\]
If a global section $s$ has zero value in a fibre of $F_*N$,
the right side pairs that zero value with $Dt$, so the left
side vanishes for every $t$. Global generation and the perfect
pairing force $Ds$ to have zero value there too.
Hence $D$ preserves every evaluation kernel and descends to an
$\mathcal O_{X^{(2)}}$-linear endomorphism of $F_*N$.
Formally the kernel of the global evaluation map is locally free;
the induced map from it to $F_*N$ vanishes on every closed fibre
and therefore vanishes identically. The descended endomorphism
acts on global sections as the original $D$.

For a line bundle $A$ of odd degree $m$ on an elliptic curve,
$F_*A$ is stable of rank two and degree $m$.
A line subbundle $T$ of degree at least $m/2$ would, by adjunction,
give a nonzero map $F^*T\to A$, hence $2\deg T\le m$,
contradicting the integral degree assumption. A noninvertible
nonzero endomorphism would have rank-one kernel and image; stability
puts both their degrees strictly below $m/2$, contradicting
additivity of degree. Every nonzero endomorphism is therefore
invertible, and over the algebraically closed field its endomorphism
algebra is $k$. Kunneth consequently gives
$\operatorname{End}(F_*N)=k$ when $m$ is odd.
Equation~\eqref{eq:ss-exact-linear-operator}, together with
$\beta(s,s)=0$, forces the scalar $D$ to be zero.

For even $m$, projection formula gives
$F_*\mathcal O(mO)=\mathcal O_{E_0^{(2)}}((m/2)O)\otimes U$,
where $U=F_*\mathcal O_{E_0}$. Its rank-one quotient in
$0\to\mathcal O\to U\to B\to0$ has degree zero. The cohomology
sequence and the zero Frobenius on $H^1(E_0,\mathcal O)$ give
$h^0(B)=1$, so $B\simeq\mathcal O$.
The extension is nonsplit because $h^0(U)=1$.
Every endomorphism preserves its unique global section and induces
scalars $a,b$ on the subline and quotient. Compatibility with its
nonzero extension class gives $a=b$. Subtracting $a$ times the
identity leaves a multiple of the quotient-followed-by-inclusion
map, which is nonzero and square zero. The invariant derivation
$\delta$ is precisely a nonzero multiple of this map. Thus
\[
 \operatorname{End}(U)=k[\delta]/(\delta^2),\qquad
 \operatorname{End}(F_*N)
       =k[\delta_1,\delta_2]/(\delta_1^2,\delta_2^2).
\]

It remains to impose the Cartier equation on
$D=a+b\delta_1+c\delta_2+e\delta_1\delta_2$.
In a pullback frame, write
$s=u^2+x_1v^2+x_2w^2+x_1x_2z^2$.
Taking the coefficient of $x_1x_2$ gives
\[
 C(s^2\omega)=C(s\delta_1s\,\omega)
 =C(s\delta_2s\,\omega)=0,
 \qquad C(s\delta_1\delta_2s\,\omega)=\delta_1\delta_2s\,\omega.
\]
Consequently $C(sD(s)\omega)=D(s)\omega$ is equivalent to
$a=b=c=0$ and $e^{1/2}=e$, or $e\in\mathbf F_2$.
The four operators are independent on global sections because
the evaluation map is surjective, so no further solutions arise.
The nonzero solution $D=\delta_1\delta_2$ exists and is nonzero:
$D(x_1x_2)=1$. This proves the geometric dimensions in both parities.

Finally $D$, the invariant differential, the line-bundle square root
and its Cartier connection are defined over $\mathbf F_2$.
The unique nonzero even-degree covector is therefore arithmetic
invariant over the stated perfect base field. In odd degree there
are no geometric invariant covectors and hence no arithmetic ones.
This proves the asserted arithmetic dimensions as well.
\end{proof}

\begin{remark}[Exact local checks and scope]
The local four-operator identities are reproduced symbolically in
\path{agents/monodromy/code/supersingular_cartier_operator_checks.py},
with output
\path{agents/monodromy/outputs/supersingular_cartier_operator_checks.json}.
The same retained program performs a bounded check of constant
alternating symmetrizers for the one-factor Cartier pencils: the
space has dimension one in pole degree three and zero in the tested
odd pole degrees five, seven, nine, eleven and thirteen. This check
is not used in the proof and asserts nothing for untested degrees;
it records why the special degree-one Hessian calculation has not
been extrapolated to higher degrees.
\end{remark}
